% Part: normal-modal-logic
% Chapter: completeness
% Section: lindenbaums-lemma

\documentclass[../../../include/open-logic-section]{subfiles}

\begin{document}

\olfileid{nml}{com}{lin}

\olsection{د لينډنباوم لمه}

د لينډنباوم لمه ثابتوي چې د !!{formula}s هر $\Sigma$-سازګار سټ
لږ تر لږه په يوه \emph{بشپړ} $\Sigma$-سازګار سټ کښې شامل دے۔
د کانوني مدل زموږ جوړونه به وښيي چې د هر بشپړ $\Sigma$-سازګار
سټ~$\Delta$ دپاره په کانوني مدل کښې داسې يوه نړۍ شته چې ټول او
يوازې هغه !!{formula}s پکې رښتيا وي چې په~$\Delta$ کښې دي۔ نو د
لينډنباوم لمه ډاډ راکوي چې هر $\Sigma$-سازګار سټ د کانوني مدل
په کومه نړۍ کښې رښتيا دے۔

\begin{thm}[د لينډنباوم لمه]\ollabel{thm:lindenbaum}
  که $\Gamma$ $\Sigma$-سازګار وي، نو داسې يو بشپړ $\Sigma$-سازګار
  سټ~$\Delta$ شته چې~$\Gamma$ غځوي۔
\end{thm}

\begin{proof}
$!A_0$، $!A_1$، \dots{} دې د ژبې د ټولو فارمولونو يو بشپړ لړ
وي (تکرارونه روا دي)۔ د بېلګې په توګه، لومړے $\Obj p_0$ په لړ
کښې راوړئ، او په هر پړاو $n \ge 1$ کښې هغه متناهي شمېر فارمولونه
په لړ کښې راوړئ چې اوږدوالے ئې تر~$n$ زيات نۀ وي او يوازې له
$\Obj p_0$، \dots،~$\Obj p_n$ څخه متغيرونه کاروي۔ په منجمده
سرچينه کښې د پړاو په دې بېلګه کښې يوازې د هماغه اوږدوالي فارمولونه
ليکل شوي دي؛ دلته د بشپړې شمېرنې دپاره تر هغه اوږدوالي پورې ټول
فارمولونه اخلو، څو د لوړو شاخصونو لنډ فارمولونه هم په لړ کښې راشي۔
د !!{formula}s سټونه~$\Delta_n$ پر~$n$ د استقرا له لارې تعريفوو،
او بيا $\Delta = \bigcup_n \Delta_n$ ږدو۔ لومړے $\Delta_0 = \Gamma$
ږدو۔ که $\Delta_n$ تعريف شوے وي، نو $\Delta_{n+1}$ داسې تعريفوو:
\[
\Delta_{n+1} =
\begin{cases}
  \Delta_n \cup \{!A_n\}, & \text{که $\Delta_n \cup \{ !A_n\}$
    $\Sigma$-سازګار وي؛} \\
  \Delta_n \cup \{ \lnot !A_n\}, & \text{که داسې نۀ وي۔}
\end{cases}
\]
اوس $\Delta = \bigcup_{n=0}^\infty \Delta_n$ ږدو۔

بايد وښيو چې دا تعريف په رښتيا داسې يو سټ~$\Delta$ راکوي چې
غوښتل شوي خاصيتونه لري، يعنې $\Gamma \subseteq \Delta$ او
$\Delta$ بشپړ $\Sigma$-سازګار دے۔

دا څرګنده ده چې $\Gamma \subseteq \Delta$، ځکه د جوړونې له مخې
$\Delta_0 \subseteq \Delta$ او $\Delta_0 = \Gamma$ دي۔ په رښتيا،
د هر~$n$ دپاره $\Delta_n \subseteq \Delta$، ځکه $\Delta$ د ټولو
سټونو~$\Delta_n$ اتحاد دے۔ (چې د جوړونې په هر ګام کښې له مخکې
جوړ شوي سټ سره !!a{formula} زياتوو، نو $\Delta_n \subseteq
\Delta_{n+1}$؛ او چې $\subseteq$ انتقالي دے، نو هر کله چې
$n \le m$ وي، $\Delta_n \subseteq \Delta_{m}$ وي۔) د جوړونې په
هر پړاو کښې يا $!A_n$ يا $\lnot !A_n$ زياتوو، او هر !!{formula}
د ټولو فارمولونو~$!A_n$ په لړ کښې (لږ تر لږه يو ځل) راځي۔ نو د
هر $!A$ دپاره يا $!A \in \Delta$ وي يا $\lnot !A \in \Delta$؛
نو د تعريف له مخې $\Delta$ بشپړ دے۔

په پای کښې بايد وښيو چې $\Delta$ $\Sigma$-سازګار دے۔ د دې دپاره
ښيو چې (a) که $\Delta$ $\Sigma$-ناسازګار وے، نو کوم $\Delta_n$
به $\Sigma$-ناسازګار وے، او (b) ټول سټونه~$\Delta_n$
$\Sigma$-سازګار دي۔

نو فرض کړئ چې $\Delta$ $\Sigma$-ناسازګار وے۔ بيا $\Delta
\Proves[\Sigma] \lfalse$، يعنې داسې $!A_1$، \dots،~$!A_k \in
\Delta$ شته چې $\Sigma \Proves !A_1 \lif (!A_2 \lif \cdots (!A_k
\lif \lfalse)\dots)$۔ چې $\Delta = \bigcup_{n=0}^\infty \Delta_n$
دے، نو هر $!A_i \in \Delta_{n_i}$ د کوم~$n_i$ دپاره دے۔ $n$ دې
په دې شاخصونو کښې تر ټولو لوے وي۔ چې $n_i \le n$، نو
$\Delta_{n_i} \subseteq \Delta_n$۔ نو ټول $!A_i$ په کوم يوه
سټ~$\Delta_n$ کښې دي۔ دا به په دې معنا وي چې $\Delta_n
\Proves[\Sigma] \lfalse$، يعنې $\Delta_n$ $\Sigma$-ناسازګار دے۔

د دې د ښودلو دپاره چې هر $\Delta_n$ $\Sigma$-سازګار دے، پر~$n$
ساده استقرا کاروو۔ $\Delta_0 = \Gamma$، او فرض مو کړے و چې
$\Gamma$ $\Sigma$-سازګار دے۔ نو دعوه د $n = 0$ دپاره سمه ده۔
اوس فرض کړئ چې د $n$ دپاره سمه وي، يعنې $\Delta_n$
$\Sigma$-سازګار وي۔ $\Delta_{n+1}$ يا $\Delta_n \cup \{!A_n\}$
دے که دا $\Sigma$-سازګار وي، او که داسې نۀ وي، نو
$\Delta_n \cup \{\lnot!A_n\}$ دے۔ په لومړي حالت کښې
$\Delta_{n+1}$ څرګند $\Sigma$-سازګار دے۔ خو د
\olref[prf][con]{prop:consistencyfacts}\olref[prf][con]{prop:consistencyfacts-c}
له مخې يا $\Delta_n \cup \{!A_n\}$ يا
$\Delta_n \cup \{\lnot!A_n\}$ سازګار دے؛ نو په دويم حالت کښې هم
$\Delta_{n+1}$ سازګار دے۔
\end{proof}

\begin{cor}\ollabel{cor:provability-characterization}
  $\Gamma \Proves[\Sigma] !A$ هله او يوازې هله وي چې $!A \in \Delta$
  د هر هغه بشپړ $\Sigma$-سازګار سټ~$\Delta$ دپاره وي چې
  $\Gamma$ غځوي (په دې کښې هغه حالت هم شامل دے چې
  $\Gamma = \emptyset$، او په هغه حالت کښې د مودالي نظام~$\Sigma$
  يوه بله ځانګړونه ترلاسه کوو)۔
\end{cor}

\begin{proof}
  فرض کړئ چې $\Gamma \Proves[\Sigma] !A$، او $\Delta$ دې هر هغه
  بشپړ $\Sigma$-سازګار سټ وي چې $\Gamma$ غځوي۔ که $!A
  \notin \Delta$ وي، نو د اعظميت له مخې $\lnot!A \in \Delta$،
  او بيا $\Delta \Proves[\Sigma] !A$ (د يکنواختۍ له مخې) او
  $\Delta \Proves[\Sigma] \lnot!A$ (د انعکاس له مخې)؛ نو
  $\Delta$ ناسازګار دے۔ په سرچپه لوري، که
  $\Gamma \Proves/[\Sigma] !A$ وي، نو
  $\Gamma \cup \{ \lnot!A\}$ $\Sigma$-سازګار دے، او د لينډنباوم
  د لمې له مخې داسې يو بشپړ سازګار سټ~$\Delta$ شته چې
  $\Gamma \cup \{ \lnot!A \}$ غځوي۔ د سازګارۍ له مخې $!A
  \notin \Delta$.
\end{proof}

\end{document}

